All terminating decimals below denote exact rationals. Here $\tau$ is the alias/checksum work parameter, $\eta n$ the large-prime bit scale, $\mu n$, $ln$ and $rn$ the sizes of the three alias blocks, $a$ the mask fraction within the large block, and $\gamma\mu n$ the small-prime bit scale. The constants $\eps$, $\zeta$ and $\theta$ specify pattern accuracy, preliminary-search slack and the bit-length guard. The frozen scalar certificate uses
\begin{center}\small
\begin{tabular}{@{}ll@{}}\toprule
Parameter & Value\\\midrule
$\tau$ & $0.48767228$\\
$\eta$ & $0.516133678389924$\\
$a$ & $0.008148495127574004$\\
$\mu$ & $0.4754356429062936$\\
$\gamma$ & $0.06005543072330252$\\
$l$ & $0.483775231690629$\\
$r=1-\mu-l$ & $0.0407891254030774$\\
$\eps,\zeta,\theta$ & $10^{-12}$ each\\
Common alias margin $s_*$ & $10^{-10}$\\
$m_{\rm fin}$ & $10^{-7}$\\
Advertised $\alpha$ & $0.4876723$\\\bottomrule
\end{tabular}
\end{center}
The exact budgets and internal total exponent are
\begin{align*}
 \Delta&=0.028461398389924,&\Gamma&=0.459210881610076,\\
 \beta=\tau+20\eps&=0.48767228002,&\alpha-\beta&=1.998\cdot10^{-8}.
\end{align*}
The construction parameters are changed, rather than only reducing the
rounding allowance in an unchanged theorem: $\tau<0.48767595$, the
internal parameter of~$\mathcal W$.

\begin{proposition}[Scalar feasibility]\label{prop:numeric}
The stated tuple satisfies every inequality in this section, with the
common alias margin $s_*=10^{-10}$, finite moment margin $10^{-7}$,
and error exponent $9.7468\cdot10^{-8}$. The continuous margin conclusion
is conditional on~\eqref{eq:oldmom}; the scalar comparisons themselves
use exact rational arithmetic with rigorous logarithm intervals.
\end{proposition}
\begin{proof}
The complete list of tested gaps and their definitions is in
Appendix~\ref{app:certificate}. In particular, all twenty alias/auxiliary
gaps exceed $s_*$, and $\zeta+\theta=2\cdot10^{-12}<s_*$.
The following decimal lower bounds have been rounded down from their
exact rational interval endpoints; their signs do not depend on a
floating-point tolerance.
\begin{center}\small
\begin{tabular}{@{}lr@{}}\toprule
Continuous margin family & Lower bound\\\midrule
Other 42 low certificates & $1.35967911\cdot10^{-7}$\\
Special filter-3 low certificate & $0.000949320934578$\\
All first bands, including extended filter 3 & $1.93876630\cdot10^{-6}$\\
Regions 2--8 & $3.8505054152$\\
Tail & $11.4443433093$\\\bottomrule
\end{tabular}
\end{center}
For clarity, the bounds are obtained from precisely these rational formulas:
\begin{align*}
 m_{\rm low}&=.00025+\eta-\eta_b-10.6428K,\\
 m_{3,0}&=.00025+\eta-\eta_b-\Lambda_0,\\
 m_{\rm first}&=.00025+\eta-\eta_b-8.4537(1.25)K,\\
 m_{2:8}&=4+\eta-\eta_b-418.36(15)K,\\
 m_9&=17+\eta-\eta_b-418.36(256)K/\Gamma_b.
\end{align*}
Every family bound exceeds $m_{\rm fin}=10^{-7}$.
Also $0<q<1$, $K>0$, $A<1/4$ and
$418.36(1+K/\Gamma_b)<419$.
Consequently the rounding-error exponent has the exact lower bound
\[
 m_{\rm fin}-(18+6\cdot419)\eps
 =10^{-7}-2532\cdot10^{-12}=9.7468\cdot10^{-8}>0.
\]
The logarithm enclosures and the exhaustive finite data checks supporting
these comparisons are proved/explained in Appendix~\ref{app:certificate}.
\end{proof}
